% ==================================================================
% id:       lim-014
% topic:    גבולות ורציפות של פונקציות
% subtopic: משפטי רציפות (ערך ביניים, ויירשטראס)
% solution: full
% ==================================================================
הוכיחו כי למשוואה
$\ds{x^4+3x=10}$
קיימים לפחות שני פתרונות בקטע $\ds{\lr{-10,10}}$.
\hint{הפעילו את משפט ערך הביניים פעמיים, על שני קטעים זרים שבכל אחד מהם הפונקציה מחליפה סימן.}
\begin{solbox}
נסמן
$$,f(x)=x^4+3x-10$$
זוהי פונקציה רציפה.
נשים לב כי
$$.f(-3)=8>0,\ f(0)=-10<0,\ f(2)=12>0$$
ע"י שימוש במשפט ערך הביניים  בקטע
$[-3,0]$
נסיק כי קיים
$x_1\in(-3,0)$
כך ש-
$f(x_1)=0$,
כלומר
$$,x_1^4+3x_1-10=0$$
זאת אומרת ש-
$x_1$
הוא פתרון של המשוואה הנתונה
(בקטע הנדון).
ע"י שימוש במשפט ערך הביניים עבור הקטע
$[0,2]$
נסיק כי קיים
$x_2\in(0,2)$
כך ש-
$f(x_2)=0$,
כלומר
$$,x_2^4+3x_2-10=0$$
זאת אומרת ש-
$x_2$
הוא פתרון של המשוואה הנתונה
(בקטע הנדון).
בבירור
$x_1\neq x_2$.
אז מצאנו למשוואה הנתונה לפחות שני פתרונות.
\end{solbox}
